
This study investigates whether Reinforcement Learning from Human Feedback (RLHF) instruction tuning degrades the discriminative quality of confidence signals in large language models. We operationalize confidence as logit margin and measure discriminative ability using AUROC for margin-based correctness prediction. Across two model families (Qwen2.5-7B and Mistral-7B), we observe that instruction-tuned models exhibit lower AUROC than their base counterparts: Qwen shows a decrease of 0.0222 (95\% CI: [0.0074, 0.0370]) and Mistral shows a decrease of 0.0385 (95\% CI: [0.0238, 0.0531]). We identify the mechanism as disproportionate margin inflation for incorrect predictions, with instruction-tuned models showing 3.06x (Qwen) and 16.79x (Mistral) higher expected margins conditional on incorrect answers compared to base models. Percentile-normalized logistic regression reveals slope attenuation of 34\% (Qwen) and 40\% (Mistral), indicating that the distortion persists after controlling for scale differences. Murphy's Brier score decomposition shows that the refinement component (discrimination) decreases in instruction-tuned models while uncertainty remains stable, consistent with geometric rather than scalar distortion. These findings suggest that RLHF may introduce a tradeoff between behavioral alignment and confidence signal reliability that standard post-hoc calibration methods targeting Expected Calibration Error may not fully address.
